% Generated draft: requires mathematical and source review.
\documentclass[11pt]{article}
\usepackage[margin=25mm]{geometry}
\usepackage{fontspec}
\setmainfont{Latin Modern Roman}
\usepackage{amsmath,amssymb,mathtools,unicode-math}
\setmathfont{Latin Modern Math}
\usepackage{enumitem,needspace,xurl,hyperref}
\setlength{\parindent}{0pt}
\setlength{\parskip}{5pt}
\setlength{\emergencystretch}{3em}
\newcommand{\N}{\mathbb{N}}
\newcommand{\Z}{\mathbb{Z}}
\newcommand{\Q}{\mathbb{Q}}
\newcommand{\R}{\mathbb{R}}
\newcommand{\C}{\mathbb{C}}
\newcommand{\F}{\mathbb{F}}
\newcommand{\A}{\mathbb{A}}
\newcommand{\PP}{\mathbb{P}}
\newcommand{\E}{\mathbb{E}}
\newcommand{\eps}{\varepsilon}
\newcommand{\dd}{\,\mathrm{d}}
\begin{document}
\section*{UnsolvedMath ID: 800559}
\textbf{Source problem code:} catalogue entry
\subsection{Definitions and mathematical statement}
The supplied source asks: ``Draw a simple polygon whose visibility graph is the given graph $G$.''

A precise candidate interpretation is the following vertex-visibility realization problem. Let $G=(V,E)$ be a finite simple undirected graph with $|V|=n\geq 3$. A simple polygon $P$ is a closed polygonal chain with distinct vertices $p_1,\ldots,p_n$, whose nonadjacent edges are disjoint and whose adjacent edges meet only at their common endpoint. Write $\overline{\Omega}_P$ for the closed region bounded by this chain, including its boundary. Under the convention that boundary visibility counts, its vertex visibility graph has vertex set $\{p_1,\ldots,p_n\}$, with distinct vertices $p_i,p_j$ adjacent exactly when
\[
[p_i,p_j]\subseteq\overline{\Omega}_P,
\]
where $[p_i,p_j]$ denotes the closed straight line segment joining them.

Under this interpretation, the task is to construct such a polygon $P$ and a bijection $f:V\to\{p_1,\ldots,p_n\}$ satisfying
\[
\forall u,v\in V,\quad u\ne v\implies
\bigl(\{u,v\}\in E\iff[f(u),f(v)]\subseteq\overline{\Omega}_P\bigr).
\]
The cyclic order of the polygon vertices is part of the desired output unless separately supplied.

This formulation is conditional on conventions absent from the supplied record: the record does not explicitly define visibility, identify the graph vertices with polygon vertices, specify degeneracy restrictions, or state whether a polygon boundary order is supplied. It also does not say whether $G$ is promised to admit a realization. Consequently, no requirement to recognize nonrealizable inputs, report nonexistence, or achieve a particular computational complexity is asserted here. These qualifications require human review against the source before this candidate formulation can be treated as definitive.
\subsection{Short English statement}
Given a graph $G$, construct a simple polygon whose vertex visibility graph is isomorphic to $G$, under the explicitly stated convention that a segment contained in the closed polygonal region gives visibility. The supplied record leaves visibility conventions, input promises, and algorithmic requirements unresolved.
\subsection{Sources}
\begin{itemize}
\item Local UnsolvedMath registry, permanent ID $800559$, exact source problem code \texttt{catalogue entry}, supplied complete source record. The registry attributes the problem to Ghosh and Goswami, with proposed year 2010, and labels it open; that label is supplied metadata, not an independently verified assessment.
\item Ghosh and Goswami, \emph{Unsolved Problems in Visibility Graphs of Points, Segments and Polygons} (2010), Problem 33. Supplied source URL: \url{https://arxiv.org/abs/1012.5187}. The record reports access on 2026-10-16 and extraction by \texttt{arxiv\_latex}. This access date and provenance are supplied information; no independent source visit or verification is claimed.
\end{itemize}
\end{document}
